\section{引言}

数学体系通常以公理化方式将“对象—关系—运算”同时置入起点。本文考虑一种互补方案：以“单一生成机制”替代多元原语，并将对象、运算与更高层概念视为该机制在有限分辨率读出下的统计稳定结构。目标并非以最少符号重述既有公理，而是提供一条严格可检验的生成链：
\begin{enumerate}
  \item 基础数据为迭代（群作用）与二元投影（$0/1$ 读出）；当讨论频率与稳定性时，引入标准的 $T$-不变概率测度并在遍历/唯一遍历条件下给出可审计结论（\cite{WaltersErgodicTheory}）；
  \item 将概念理解为“可寻址事件”，地址来自旧概念读出的有限序列；
  \item 以统计稳定性（极限频率/不变状态）将可忽略噪声与不可忽略结构区分；
  \item 以黄金分支 $\varphi^{-1}$ 作为\textbf{在 Hurwitz/Markov（最坏有理逼近/反共振）口径下的规范代表}，提供自相似重整化的最简示例、最坏情形的反共振界限与与 Zeckendorf/Ostrowski 语法一致的递归编码（更一般“稳定性指标”下的唯一性不在本文主张范围内）；
  \item 证明自然数、加法、乘法可作为稳定地址上的归一化运算涌现。
\end{enumerate}

本文强调两点：
\begin{itemize}
  \item “从无到有”在此并非指绝对空集意义的无，而是指\textbf{在不给定地址之前概念值不定义}；只有当地址被旧概念读出所指派，概念值才作为二元投影出现。
  \item “运算涌现”在此并非心理学或经验主义陈述，而是指：在稳定地址空间上存在\textbf{可计算、可终止}的规范化过程（例如 Zeckendorf 贪心规范化），其诱导运算与经典算术同构。
\end{itemize}

\subsection{阅读指南：数学结果与解释性框架的分层}\label{subsec:reading-guide}
本手稿包含两类文本，它们在逻辑地位上严格不同：
\begin{itemize}
  \item \textbf{可证结果（以定义/引理/定理/命题呈现）}：这些结论只依赖明确给出的数学对象与引用的标准定理；任何后续使用都以编号引用为准。
  \item \textbf{解释性叙事（以 remark/讨论段落呈现）}：例如“概念/运算的涌现”等表述，仅用于动机与直观解释，不作为后续证明的前提。为避免遮蔽数学贡献，本文尽量将这类语言放入注（remark）或结论讨论，而不混入定理陈述。
\end{itemize}

\subsection{贡献清单与非主张（scope control）}\label{subsec:contrib-scope}
\textbf{本文的数学贡献}可概括为以下几条“可审计接口”：
\begin{itemize}
  \item \textbf{反共振极值口径：}给出 Hurwitz/Markov 意义下黄金分支的极值与（按模群等价的）唯一性，并将其转译为差异度证书的可比较上界（第 2 节与附录 H）。
  \item \textbf{分辨率折叠与多尺度对象：}构造可计算折叠 $\Fold_m$ 并证明其基本性质与跨尺度相容；用逆系统/逆极限把“不同 $m$ 仅是可见前缀长度不同”严格化（第 6 节与附录 G）。
  \item \textbf{序列级因子与有限状态表示：}把窗口层折叠提升为滑动块算子 $\Phi_m$，给出因子性质与显式 sofic 图表示（第 6 节）。
  \item \textbf{算术与归一化：}把 Zeckendorf 规范代表视为值纤维的横截面，并给出“逐位叠加 $\Rightarrow$ 折叠归一化”的等价表述；同时引用 Zeckendorf 数制上纯数字层（有限状态/线性时间）的加法实现（第 7 节与附录 E；参见 \cite{Frougny1992NumbersAutomata,AhlbachUsatineFrougnyPippenger2013}）。
  \item \textbf{核对比：}把黄金比例基上 $9/13/21$-local 的三类并行归一化核统一为“零式 $+$ $q$-场”模板，并在单流在线机层面给出可 BFS 的最小状态规格、carry-free 骨架的 $\zeta$ 三分分类与 primitive（“素数影子”）统一提取接口（第 \ref{sec:A-kernel-compare-main} 小节与附录 \ref{app:sync-kernel-A-compare}）。
  \item \textbf{Parry/\(\zeta\)/算子接口：}在黄金均值 SFT 上给出 Parry 最大熵测度与动力学 $\zeta$ 的标准闭式表达，并将柱投影/对角子代数/态的口径写成可核验命题（第 8--9 节）。
\end{itemize}

\textbf{本文明确不主张}下列更强陈述（除非额外假设另行给出）：
\begin{itemize}
  \item \textbf{旋转（Sturmian）下折叠极限等于 Parry：}一般不成立；本文将 Parry 分布作为最大熵基线，并用误差分解定位“有限样本误差”与“模型差异项”（第 \ref{sec:experiments} 节）。
  \item \textbf{哲学性用语具有证明效力：}诸如“概念涌现/时间涌现”等仅作解释性语言，不参与推导链（见小节 \ref{subsec:reading-guide}）。
\end{itemize}

\subsection{本文的可检验输出（定量）}
本文的主要输出均可在纯数学意义下被审计：
\begin{itemize}
  \item \textbf{计数律：}稳定类型空间 $X_m$ 的基数为 $|X_m|=F_{m+2}$，并给出折叠压缩率 $2^m/F_{m+2}$（第 \ref{sec:folding} 节与附录 B--D）。
  \item \textbf{极限频率：}在唯一遍历/遍历条件下，稳定事件的轨道频率收敛到不变测度（第 4 节，参见 \cite{WaltersErgodicTheory}）。
  \item \textbf{复杂度证书：}在圆周无理旋转的二分割读出模型中，输出序列为 Sturmian，因子复杂度满足 $p(n)=n+1$（定理 \ref{thm:sturmian-min-complexity-rotation}）。
  \item \textbf{拟合准则（无需模拟）：}对任何观测到的稳定类型序列，可用经验转移频率与 Perron--Frobenius 平稳分布比较，给出总变差/相对熵等偏差指标（第 8 节，参见 \cite{LindMarcus1995}）。
  \item \textbf{算术同构：}以 Zeckendorf 规范形为稳定地址，定义运算并由值映射与唯一性得到与 $(\NN,+,\times)$ 的同构（第 7 节，参见 \cite{Zeckendorf1972}）。
  \item \textbf{解析接口：}对黄金均值约束（有限型子移位）给出动力学 $\zeta$ 的闭式表达（第 9 节，参见 \cite{LindMarcus1995}）。
\end{itemize}

\subsection{与既有结果的关系（仅列引用点）}
本文复用三条成熟工具链：符号动力系统与有限型子移位 \cite{LindMarcus1995}、遍历理论的频率极限定理 \cite{WaltersErgodicTheory}、以及 Zeckendorf 表示的存在唯一性 \cite{Zeckendorf1972}。黄金比例的“最难有理逼近”性质将作为反共振选择的严格形式（第 2 节，参见 \cite{Cassels1957,KuipersNiederreiter1974}）。

